% Part: incompleteness
% Chapter: introduction
% Section: overview

\documentclass[../../../include/open-logic-section]{subfiles}

\begin{document}

\olfileid[id]{inc}{int}{ovr}

\olsection{Ikhtisar Hasil-Hasil Ketaklengkapan}

Hilbert berharap bahwa matematika dapat diformalkan dalam suatu teori
!!{axiomatizable} yang dapat dibuktikan !!{complete} dan
!!{decidable}. Selain itu, ia bermaksud membuktikan konsistensi teori
tersebut dengan sarana yang sangat lemah dan ``finitistik,'' sehingga matematika
klasik dapat dipertahankan dari tantangan intuisionisme. Teorema-teorema
ketaklengkapan G\"odel menunjukkan bahwa tujuan-tujuan ini tidak dapat dicapai.

Teorema ketaklengkapan pertama G\"odel menunjukkan bahwa suatu versi
\emph{Principia Mathematica} karya Russell dan Whitehead tidak
!!{complete}. Namun, pembuktiannya sebenarnya sangat umum dan berlaku bagi
berbagai macam teori. Artinya, persoalannya bukan sekadar bahwa
\emph{Principia Mathematica} gagal menangkap matematika secara lengkap,
melainkan bahwa \emph{tidak ada} teori yang memadai yang dapat melakukannya.
Diperlukan waktu untuk mengisolasi ciri-ciri teori yang cukup agar
teorema-teorema ketaklengkapan berlaku, lalu menggeneralisasi pembuktian
G\"odel sehingga hanya bergantung pada ciri-ciri tersebut. Kini kita dapat
menyatakan versi yang sangat umum dari teorema ketaklengkapan pertama bagi
teori-teori dalam bahasa aritmetika~$\Lang L_A$.

\begin{thm}
Jika $\Gamma$ adalah teori yang konsisten dan !!{axiomatizable} dalam~$\Lang
L_A$ yang !!{represents} semua fungsi yang dapat dihitung dan relasi
!!{decidable}, maka $\Gamma$ tidak !!{complete}.
\end{thm}

Pernyataan bahwa $\Gamma$ tidak !!{complete} berarti bahwa terdapat
sekurang-kurangnya satu !!{sentence}~$!A$ sedemikian sehingga $\Gamma
\Proves/ !A$ dan $\Gamma \Proves/ \lnot !A$. !!a{sentence} semacam itu
disebut \emph{independen} (terhadap~$\Gamma)$. Kita sebenarnya dapat
membuktikan dengan relatif cepat bahwa kalimat independen pasti ada. Namun,
kekuatan pembuktian G\"odel terletak pada kenyataan bahwa pembuktian itu
menampilkan sebuah \emph{contoh khusus} dari !!{sentence} independen
tersebut. Konstruksi yang menarik ini menghasilkan !!a{sentence}~$!G_\Gamma$,
yang disebut \emph{kalimat G\"odel} bagi~$\Gamma$, yang tidak dapat dibuktikan
karena di dalam $\Gamma$, $!G_\Gamma$ ekuivalen dengan klaim bahwa
$!G_\Gamma$ tidak dapat dibuktikan dalam~$\Gamma$. Konstruksi tersebut
bersifat \emph{konstruktif}: jika diberikan aksiomatisasi~$\Gamma$ dan
deskripsi sistem !!{derivation}, pembuktiannya memberikan metode untuk
benar-benar menuliskan~$!G_\Gamma$.

Konstruksi dalam pembuktian G\"odel mengharuskan kita menemukan cara untuk
menyatakan dalam $\Lang L_A$ sifat dan operasi pada suku serta !!{formula}s
dari $\Lang L_A$ itu sendiri. Di antaranya ialah sifat seperti ``$!A$ adalah
!!a{sentence},'' ``$\delta$ adalah !!a{derivation} dari~$!A$,'' dan operasi
seperti $\Subst{!A}{t}{x}$. Cara tersebut harus (a) menyatakan sifat dan
relasi ini melalui suatu ``pengodean'' simbol dan barisannya (karena suku,
!!{formula}s, !!{derivation}s, dan sebagainya merupakan barisan simbol) sebagai
bilangan asli (yang dapat dibicarakan oleh $\Lang L_A$). Cara itu juga harus
(b) melakukannya sedemikian rupa sehingga $\Gamma$ membuktikan fakta-fakta
yang relevan; jadi kita harus menunjukkan bahwa sifat-sifat tersebut dikodekan
oleh sifat !!{decidable} dari bilangan asli dan bahwa operasinya bersesuaian
dengan fungsi yang dapat dihitung pada bilangan asli. Proses ini disebut
``aritmetisasi sintaksis.''

Namun, sebelum menyelidiki cara mengaritmetisasi sintaksis, kita akan
mempertimbangkan syarat bahwa~$\Gamma$ ``cukup kuat,'' yakni !!{represents}
semua fungsi yang dapat dihitung dan relasi !!{decidable}. Hal ini mengharuskan
kita memberikan definisi yang tepat bagi ``dapat dihitung.'' Definisi tersebut
dapat diberikan dengan berbagai cara, misalnya melalui model mesin Turing
atau sebagai fungsi yang dapat dihitung oleh program dalam suatu bahasa
pemrograman serbaguna. Akan tetapi, karena tujuan kita ialah
!!{represents}s fungsi dan relasi tersebut dalam suatu teori berbahasa~$\Lang
L_A$, sebaiknya kita memilih definisi sederhana tentang keterhitungan fungsi
numerik saja. Definisi itu ialah gagasan \emph{fungsi rekursif}. Karena itu,
mula-mula kita akan membahas fungsi rekursif. Selanjutnya kita akan menunjukkan
bahwa $\Th{Q}$ sudah !!{represents} semua fungsi dan relasi rekursif. Hasil ini
memungkinkan kita menerapkan teorema ketaklengkapan pada teori tertentu seperti
$\Th{Q}$ dan $\Th{PA}$, sebab kita telah menetapkan bahwa keduanya merupakan
contoh teori yang ``cukup kuat.''

Hasil akhir aritmetisasi sintaksis ialah !!a{formula} $\OProv[\Gamma](x)$
yang, melalui pengodean !!{formula}s sebagai bilangan, menyatakan
keterbuktian dari aksioma-aksioma~$\Gamma$. Secara khusus, jika $!A$ dikodekan
oleh bilangan~$n$ dan $\Gamma \Proves !A$, maka $\Gamma \Proves
\OProv[\Gamma](\num{n})$. ``Predikat keterbuktian'' bagi $\Gamma$ ini juga
memungkinkan kita menyatakan, dalam arti tertentu, konsistensi $\Gamma$
sebagai !!a{sentence} dalam~$\Lang L_A$: tetapkan ``pernyataan konsistensi''
bagi~$\Gamma$ sebagai !!{sentence} $\lnot \OProv[\Gamma](\num{n})$, dengan
$n$ sebagai kode suatu kontradiksi, misalnya~$\lfalse$. Teorema
ketaklengkapan kedua menyatakan bahwa teori konsisten yang !!{axiomatizable}
juga tidak membuktikan pernyataan konsistensinya sendiri. Syarat agar teorema
ini berlaku sedikit lebih ketat daripada sekadar bahwa teori tersebut
merepresentasikan semua fungsi yang dapat dihitung dan relasi !!{decidable},
tetapi kita akan menunjukkan bahwa $\Th{PA}$ memenuhinya.

\end{document}
